% Generated draft: requires mathematical and source review.
\documentclass[11pt]{article}
\usepackage[margin=25mm]{geometry}
\usepackage{fontspec}
\setmainfont{Latin Modern Roman}
\usepackage{amsmath,amssymb,mathtools,unicode-math}
\setmathfont{Latin Modern Math}
\usepackage{enumitem,needspace,xurl,hyperref}
\setlength{\parindent}{0pt}
\setlength{\parskip}{5pt}
\setlength{\emergencystretch}{3em}
\newcommand{\N}{\mathbb{N}}
\newcommand{\Z}{\mathbb{Z}}
\newcommand{\Q}{\mathbb{Q}}
\newcommand{\R}{\mathbb{R}}
\newcommand{\C}{\mathbb{C}}
\newcommand{\F}{\mathbb{F}}
\newcommand{\A}{\mathbb{A}}
\newcommand{\PP}{\mathbb{P}}
\newcommand{\E}{\mathbb{E}}
\newcommand{\eps}{\varepsilon}
\newcommand{\dd}{\,\mathrm{d}}
\begin{document}
\section*{UnsolvedMath ID: 801159}
\textbf{Source problem code:} catalogue entry
\subsection{Definitions and mathematical statement}
A standard underlying Frobenius structure consists of a manifold \(M\), a nondegenerate symmetric metric \(g\), a multiplication \(\circ\) on \(TM\), and a unit vector field \(e\). The multiplication is bilinear, commutative and associative on each tangent space, and
\[
 g(X\circ Y,Z)=g(X,Y\circ Z),\qquad e\circ X=X.
\]
Writing \(\nabla\) for the Levi--Civita connection of \(g\) and \(c(X,Y,Z)=g(X\circ Y,Z)\), the usual flat formulation requires
\[
 R^\nabla=0,\qquad \nabla e=0,\qquad \nabla c\text{ is totally symmetric}.
\]
The source does not specify the real or complex category, regularity, dimension, metric signature, or whether an Euler vector field and its homogeneity conditions are part of the intended Frobenius structure. The displayed axioms provide terminology rather than resolve these choices.

The supplied Problem 12 asks for clarification and construction in the following directions.
\begin{enumerate}
\item Is Frobenius structure projectively invariant? One standard meaning of projective equivalence of torsion-free connections \(\nabla\) and \(\widehat\nabla\) is the existence of a one-form \(\alpha\) such that
\[
 \widehat\nabla_XY=\nabla_XY+\alpha(X)Y+\alpha(Y)X
\]
for all vector fields \(X,Y\). The source does not specify whether this is its intended meaning, which projective transformations are allowed, or how the metric, multiplication, unit, and any Euler field should transform. Consequently, an exact invariance assertion cannot be extracted from the supplied record.
\item Can the flat metric in the definition of Frobenius structure be replaced by a metric of constant sectional curvature? The source does not specify which remaining Frobenius axioms must be retained or modified, whether nonzero curvature is required, or whether the target is a consistent definition, examples, or a general construction.
\item Is there a construction assigning a Killing tensor to a Frobenius manifold? For terminology, a symmetric covariant tensor \(K\) of rank \(r\geq 1\) is a Killing tensor for a metric \(h\) when its Levi--Civita connection \(\nabla^h\) satisfies
\[
 \nabla^h_{(i_0}K_{i_1\ldots i_r)}=0,
\]
where parentheses denote symmetrization. The source motivates this question by saying that certain Frobenius manifolds have been used to construct Killing tensors associated with geometric Poisson brackets of order \(3\). It further describes these brackets as essentially equivalent to Killing tensors satisfying an additional relation on their covariant derivative. That relation, the tensor rank, the background metric, and the precise class of geometric Poisson brackets are absent from the supplied record. The intended construction therefore needs these specifications; the ordinary Killing equation alone does not encode the stated target.
\end{enumerate}
The source cautiously suggests that certain Frobenius manifolds give examples of geometric Poisson brackets of order \(3\), adding that the WDVV equation does so. This is retained as supplied motivation, without asserting a construction for every Frobenius manifold. The scope of the last question---all Frobenius manifolds or a specified subclass---is unresolved.
\subsection{Short English statement}
Problem 12 asks whether Frobenius structures have a projective invariance, whether their flat metric can be replaced by a constant-curvature metric, and whether they produce the special Killing tensors associated with third-order geometric Poisson brackets. The supplied record leaves the transformation rules, modified axioms, and additional tensor equation unspecified. Its status is supplied as open; this draft does not independently verify that status.
\subsection{Sources}
\begin{itemize}
\item Local UnsolvedMath registry, permanent ID \(801159\), exact source problem code \texttt{catalogue entry}, title ``Nijenhuis Geometry and Integrable Systems --- Problem 12.'' Supplied upstream revision: \texttt{d371ee917ac9647fbb06cd6fc529deec6f3dda82}. The registry labels the problem open; this is a supplied status, not an independently verified mathematical assessment.
\item Bolsinov, Konyaev, Matveev, \emph{Research Problems on Nijenhuis Geometry and Integrable Systems} (2024), Problem 12, as attributed in the supplied record: \url{https://arxiv.org/abs/2410.04276}. The registry also attributes the proposal to these authors in 2024.
\item Supplied provenance records access on 2026-10-16 and extraction by \texttt{arxiv\_latex}. These are supplied provenance details; the source has not been independently visited or verified for this draft. The supplied passage refers to preceding discussion for the additional covariant-derivative relation, but that discussion is not included.
\end{itemize}
\end{document}
